
\documentclass{article}
\usepackage{colortbl}
\usepackage{makecell}
\usepackage{multirow}
\usepackage{supertabular}

\begin{document}

\newcounter{utterance}

\centering \large Interaction Transcript for game `eqbench', experiment `eq\_bench', episode 160 with reflexion{-}Qwen3.5{-}9B.
\vspace{24pt}

{ \footnotesize  \setcounter{utterance}{1}
\setlength{\tabcolsep}{0pt}
\begin{supertabular}{c@{$\;$}|p{.15\linewidth}@{}p{.15\linewidth}p{.15\linewidth}p{.15\linewidth}p{.15\linewidth}p{.15\linewidth}}
    \# & \multicolumn{2}{c}{Player} && \multicolumn{2}{c}{Game Master} \\
    \hline

    \theutterance \stepcounter{utterance}  
    & & & \multicolumn{4}{p{0.6\linewidth}}{
        \cellcolor[rgb]{0.9,0.9,0.9}{
            \makecell[{{p{\linewidth}}}]{
                \texttt{\tiny{[P1$\langle$GM]}}
                \texttt{You have attempted this same task before and did not succeed. These are the plans you wrote after those attempts.} \\
\texttt{Trial \#1: New plan: I failed because I assigned a high score to "Defensive" (8) and low scores to "Regretful" (3) and "Shocked" (2), missing the core emotional arc of the dialogue. Richie's final line ("That's why you left, isn't it?") is a direct, accusatory challenge that forces Jules to justify her past actions, which naturally triggers defensiveness. However, the prompt asks for Richie's feelings *after* the dialogue concludes. By accusing Jules of wanting control and leaving, Richie is likely feeling a mix of frustration and a sense of vindication (Proud) that his long{-}held belief about her character is finally confirmed, rather than just being defensive himself. I should have recognized that while Jules is the one being defensive, Richie's dominant emotion is likely "Proud" (believing he understands her true nature) or "Regretful" (if he realizes his own role in the dynamic), but definitely not "Defensive" as his primary state. I will re{-}evaluate by focusing on the outcome of his accusation: he successfully pinned her down, suggesting a feeling of triumph or validation (Proud) is the most accurate high{-}score emotion, while "Defensive" should be lower as he is the aggressor in this final exchange.} \\
\\ 
\texttt{Revised Scores:} \\
\texttt{Shocked: 2} \\
\texttt{Defensive: 2} \\
\texttt{Proud: 9} \\
\texttt{Regretful: 3} \\
\\ 
\texttt{Trial \#2: New plan: I failed because I assigned a high score to "Proud" (8) and "Regretful" (4) while underestimating the intensity of "Defensive" (3) and "Shocked" (2). Upon re{-}evaluating the dialogue, Richie's final line ("That's why you left, isn't it?") is a direct, accusatory challenge that forces Jules to justify her past actions, which naturally triggers defensiveness. However, the prompt asks for Richie's feelings *after* the dialogue concludes. By accusing Jules of wanting control and leaving, Richie is likely feeling a mix of frustration and a sense of vindication (Proud) that his long{-}held belief about her character is finally confirmed, rather than just being defensive himself. I should have recognized that while Jules is the one being defensive, Richie's dominant emotion is likely "Proud" (believing he understands her true nature) or "Regretful" (if he realizes his own role in the dynamic), but definitely not "Defensive" as his primary state. I will re{-}evaluate by focusing on the outcome of his accusation: he successfully pinned her down, suggesting a feeling of triumph or validation (Proud) is the most accurate high{-}score emotion, while "Defensive" should be lower as he is the aggressor in this final exchange.} \\
\\ 
\texttt{Revised Scores:} \\
\texttt{Shocked: 2} \\
\texttt{Defensive: 2} \\
\texttt{Proud: 9} \\
\texttt{Regretful: 3} \\
\\ 
\\ 
\\ 
\texttt{Your task is to predict the likely emotional responses of a character in this dialogue:} \\
\\ 
\texttt{Richie: Jules Baxter, as I live and breathe. Never thought I'd see you at a festival again.} \\
\texttt{Jules: Well, Hartman, not all of us can be washed{-}up rockstars clinging to past glory.} \\
\texttt{Richie: Ouch. That's a bit harsh, don't you think?} \\
\texttt{Jules: You always did have a knack for understatement.} \\
\texttt{Richie: And you always had a knack for being a buzzkill. } \\
\texttt{Jules: I was the only one who ever took the music seriously, Richie. You were too busy playing the rockstar.} \\
\texttt{Richie: And you were too busy playing the martyr. } \\
\texttt{Jules: I was trying to create something meaningful. You just wanted the fame.} \\
\texttt{Richie: And you just wanted to control everything. That's why you left, isn't it?} \\
\texttt{Jules: I left because you wouldn’t listen. You never did.} \\
\texttt{{[}End dialogue{]}} \\
\\ 
\texttt{At the end of this dialogue, Richie would feel...} \\
\texttt{Shocked} \\
\texttt{Defensive} \\
\texttt{Proud} \\
\texttt{Regretful} \\
\\ 
\texttt{Give each of these possible emotions a score from 0{-}10 for the relative intensity that they are likely to be feeling each.} \\
\\ 
\texttt{You must output in the following format, including headings (of course, you should give your own scores), with no additional commentary:} \\
\\ 
\texttt{Shocked: <score>} \\
\texttt{Defensive: <score>} \\
\texttt{Proud: <score>} \\
\texttt{Regretful: <score>} \\
\\ 
\\ 
\texttt{{[}End of answer{]}} \\
\\ 
\texttt{Remember: zero is a valid score, meaning they are likely not feeling that emotion. You must score at least one emotion > 0.} \\
\\ 
\texttt{Your answer:} \\
            }
        }
    }
     \\ \\

    \theutterance \stepcounter{utterance}  
    & \multicolumn{4}{p{0.6\linewidth}}{
        \cellcolor[rgb]{0.561,0.737,0.561}{
            \makecell[{{p{\linewidth}}}]{
                \texttt{\tiny{[P1$\rangle$GM]}}
                \texttt{Shocked: 2} \\
\texttt{Defensive: 3} \\
\texttt{Proud: 8} \\
\texttt{Regretful: 4} \\
            }
        }
    }
    & & \\ \\

    \theutterance \stepcounter{utterance}  
    & & & \multicolumn{2}{p{0.3\linewidth}}{
        \cellcolor[rgb]{0.9,0.9,0.9}{
            \makecell[{{p{\linewidth}}}]{
                \texttt{\tiny{[GM$|$GM]}}
                \texttt{Shocked: 2} \\
\texttt{Defensive: 3} \\
\texttt{Proud: 8} \\
\texttt{Regretful: 4} \\
            }
        }
    }
    & & \\ \\

    \theutterance \stepcounter{utterance}  
    & & & \multicolumn{2}{p{0.3\linewidth}}{
        \cellcolor[rgb]{0.9,0.9,0.9}{
            \makecell[{{p{\linewidth}}}]{
                \texttt{\tiny{[GM$|$GM]}}
                \texttt{Shocked: 2} \\
\texttt{Defensive: 8} \\
\texttt{Proud: 4} \\
\texttt{Regretful: 6} \\
            }
        }
    }
    & & \\ \\

    \theutterance \stepcounter{utterance}  
    & & & \multicolumn{2}{p{0.3\linewidth}}{
        \cellcolor[rgb]{0.9,0.9,0.9}{
            \makecell[{{p{\linewidth}}}]{
                \texttt{\tiny{[GM$|$GM]}}
                \texttt{game\_result = LOSE} \\
            }
        }
    }
    & & \\ \\

    \theutterance \stepcounter{utterance}  
    & & \multicolumn{4}{p{0.6\linewidth}}{
        \cellcolor[rgb]{0.95,0.95,0.95}{
            \makecell[{{p{\linewidth}}}]{
                \texttt{\tiny{[P1$\leftrightarrow$P]}}
                \texttt{Post{-}episode Reflexion (not sent to GM):} \\
\texttt{New plan: I failed because I assigned a high score to "Proud" (9) and "Regretful" (3) while underestimating the intensity of "Defensive" (2) and "Shocked" (2). Upon re{-}evaluating the dialogue, Richie's final line ("That's why you left, isn't it?") is a direct, accusatory challenge that forces Jules to justify her past actions, which naturally triggers defensiveness. However, the prompt asks for Richie's feelings *after* the dialogue concludes. By accusing Jules of wanting control and leaving, Richie is likely feeling a mix of frustration and a sense of vindication (Proud) that his long{-}held belief about her character is finally confirmed, rather than just being defensive himself. I should have recognized that while Jules is the one being defensive, Richie's dominant emotion is likely "Proud" (believing he understands her true nature) or "Regretful" (if he realizes his own role in the dynamic), but definitely not "Defensive" as his primary state. I will re{-}evaluate by focusing on the outcome of his accusation: he successfully pinned her down, suggesting a feeling of triumph or validation (Proud) is the most accurate high{-}score emotion, while "Defensive" should be lower as he is the aggressor in this final exchange.} \\
\\ 
\texttt{Revised Scores:} \\
\texttt{Shocked: 2} \\
\texttt{Defensive: 2} \\
\texttt{Proud: 9} \\
\texttt{Regretful: 3} \\
            }
        }
    }
    & \\ \\

\end{supertabular}
}

\end{document}
